% !TeX root = ../../thesis.tex

\chapter{Alignment and Updating of the Data}
\label{ch:4_Alignment_Updating}

% ── New glossary entries ─────────────────────────────────────────────────────
\newglossaryentry{GNSS}{name={GNSS}, description={Global Navigation Satellite System}}
\newglossaryentry{IMU}{name={IMU}, description={Inertial Measurement Unit}}
\newglossaryentry{BLE}{name={BLE}, description={Bluetooth Low Energy}}
\newglossaryentry{SfM}{name={SfM}, description={Structure-from-Motion}}
\newglossaryentry{ORB}{name={ORB}, description={Oriented FAST and Rotated BRIEF}}
\newglossaryentry{FPFH}{name={FPFH}, description={Fast Point Feature Histograms}}
\newglossaryentry{RMSE}{name={RMSE}, description={Root Mean Square Error}}
\newglossaryentry{RDF}{name={RDF}, description={Resource Description Framework}}
\newglossaryentry{SUPER4PCS}{name={SUPER4PCS}, description={Super 4-Points Congruent Sets}}
\newglossaryentry{ICP}{name={ICP}, description={Iterative Closest Point}}
\newglossaryentry{API}{name={API}, description={Application Programming Interface}}
\newglossaryentry{LCP}{name={LCP}, description={Largest Common Pointset}}

%- Paper references for this chapter
\definecolor{shadecolor}{gray}{0.9}
\markerlessfootnote{
	\begin{shaded}
		Parts of this chapter's content were previously published in \cite{vermandere_automatic_2022,vermandere_two-step_2022}.
	\end{shaded}
}

% ─────────────────────────────────────────────────────────────────────────────
\section{Introduction}
\label{sec:4_intro}

% Paragraph purpose: situate the chapter in the pipeline and frame the core technical problem
This chapter addresses the first stage of the \gls{DRM} pipeline: aligning newly captured data with pre-existing reference datasets and incrementally updating those datasets with new observations (Figure~\ref{fig:4_align_update_intro}). As established in Chapter~\ref{ch:introduction}, the absence of \gls{GNSS} in indoor environments means that the pose of a newly captured session cannot be determined directly, but must instead be estimated relative to existing reference data~\cite{vermandere_two-step_2022}. This is compounded by temporal change: objects may have moved or been replaced since the reference was acquired, causing feature-based co-registration methods that assume a static scene to fail~\cite{vermandere_two-step_2022}, while full re-documentation with high-end equipment remains impractical~\cite{bassier_standalone_2016, bassier_geomapi_2024}.

\begin{figure}[h]
  \centering
  \includegraphics[width=1\textwidth]{Align_update_intro}
  \caption[Alignment and Updating]{A reference scene (left) and a new session (middle) are aligned to produce a combined dataset (right) that integrates the new observations while preserving unchanged regions of the original reference.}
  \label{fig:4_align_update_intro}
\end{figure}

% Paragraph purpose: introduce the proposed solution
To address both challenges, this chapter presents a two-step alignment framework combined with a dataset updating mechanism. The alignment combines a coarse global positioning step with a fine multi-modal registration step that jointly leverages 2D image features and 3D geometric descriptors to produce a robust pose estimate. The updating procedure then determines which points in the existing dataset remain current and which should be replaced by the new observations.

% ─────────────────────────────────────────────────────────────────────────────
\section{Related Work}
\label{sec:4_related}

This section reviews the four areas of prior work most relevant to the contributions of this chapter: Low-Cost Sensor Alignment, 3D-Based Alignment, 2D-Based Alignment and Temporal Updating.

% ─────────────────────────────────────────────────────────────────────────────
\subsection{Low-Cost Sensor Data Alignment}
\label{sec:4_related_lowcost}

% Paragraph purpose: introduce the challenge of aligning low-cost sensor data
The alignment of data captured by low-cost sensors such as \gls{XR} devices and smartphones presents unique challenges compared to the registration of high-end scanner data. These devices produce measurements of significantly lower accuracy, density, and spatial extent, which limits the effectiveness of conventional registration algorithms~\cite{geyter_point_2022}.

% Paragraph purpose: discuss local tracking and SLAM
Within a single measurement session, low-cost devices maintain their local pose through \gls{SLAM} algorithms that combine visual and inertial measurements. For instance, the Microsoft HoloLens~2 employs ORB-SLAM~\cite{cyrus_hololens_2019} using four grey-scale tracking cameras and an \gls{IMU} to continuously estimate the device's position and orientation within a local coordinate system. Although drift accumulates significantly over longer multi-room trajectories~\cite{geyter_point_2022, chen_accuracy_2018}, these local tracking methods are sufficiently accurate for single-room scenes, achieving accuracies conforming to LOA20 ($2\sigma \leq 0.05$\,m)~\cite{geyter_point_2022}. However, the resulting data remains in an arbitrary local coordinate frame that bears no relation to any global or facility-wide reference system.

% Paragraph purpose: discuss the gap between local tracking and global registration
Bridging this gap between accurate local tracking and global positioning is the central challenge addressed in this chapter. The local SLAM provides reliable relative positioning within a session, but the entire session must still be placed within the coordinate system of the pre-existing reference data. This is fundamentally a one-time transformation estimation per session, rather than a continuous tracking problem.

% ─────────────────────────────────────────────────────────────────────────────
\subsection{3D-Based Alignment}
\label{sec:4_related_3d}

% Paragraph purpose: review geometric registration approaches
Three-dimensional registration methods estimate the transformation between two geometric datasets by matching corresponding features in 3D space. The most established approach is the Iterative Closest Point (\gls{ICP}) algorithm and its variants, which iteratively minimise the distance between corresponding point pairs. While effective for fine alignment, ICP requires a reasonable initial estimate and is sensitive to partial overlap and noise~\cite{bassier_geomapi_2024}.

\begin{wrapfigure}{r}{0.4\textwidth}
  \centering
  \includegraphics[width=0.4\textwidth]{BIM_Feature_example}
  \caption[BIM features]{Example of limited BIM features (green) that make traditional alignmet approaches unreliable.}
  \label{fig:BIM_Features}
\end{wrapfigure}
% Paragraph purpose: discuss feature-based global registration
For coarse alignment without a prior estimate, feature-based methods extract distinctive geometric descriptors from each dataset and establish correspondences between them. \gls{FPFH} descriptors~\cite{rusu_fast_2009} encode the local surface geometry around each point through histograms of angular relationships between surface normals. These descriptors can be matched across datasets to establish candidate correspondences, from which a rigid body transformation is estimated. The Fast Global Registration method proposed by Zhou et al.~\cite{zhou_fast_2016} avoids computationally demanding RANSAC-based outlier rejection by instead using a robust correspondence estimation function that simultaneously minimises distances between correspondences and down-weights outliers.

% Paragraph purpose: discuss BIM-based alignment
When a \gls{BIM} model is available, it can serve as an additional reference for geometric alignment. However, BIM geometries are typically abstract and sparse compared to captured point clouds, making conventional point-based features unreliable (Figure~\ref{fig:BIM_Features}). The \gls{SUPER4PCS} algorithm~\cite{mellado_super_2014} addresses this by matching approximately congruent four-point sets between datasets, relying on geometric ratios rather than local surface detail. This makes it particularly well suited for aligning captured data with the planar and simplified geometry of BIM models. A key computational improvement over the original 4PCS algorithm is the use of a smart indexing structure that retrieves, for a given pair of points, all other pairs at approximately the same distance directly, without testing every combination. This reduces the pair-extraction cost from $O(n^2)$ to an expected $O(n)$, lowering the overall complexity from $O(n^2+k)$ to $O(n+k)$, where $k$ is the number of reported congruent sets.

% ─────────────────────────────────────────────────────────────────────────────
\subsection{2D-Based Alignment}
\label{sec:4_related_2d}

% Paragraph purpose: review image-based registration approaches
Image-based alignment methods estimate the pose of a sensor by matching visual features between newly captured images and a set of geo-referenced reference images. This approach is particularly relevant for \gls{XR} devices and smartphones, which are equipped with cameras that produce imagery alongside their geometric measurements.

% Paragraph purpose: discuss feature matching for pose estimation
Structure-from-Motion (\gls{SfM}) techniques estimate camera poses by detecting and matching distinctive image features across multiple views~\cite{schonberger_structure--motion_2016}. A wide range of local feature detectors and descriptors exists for this purpose, including SIFT~\cite{sun_sift_2014}, SURF~\cite{bay_speeded-up_2008}, and \gls{ORB}~\cite{rublee_orb_2011}, each offering different trade-offs between distinctiveness and computational cost. \gls{ORB} features in particular provide a computationally efficient basis for this matching and are used, for instance, by the Microsoft HoloLens~2~\cite{cyrus_hololens_2019}. More recently, learned alternatives have replaced these handcrafted descriptors with features trained directly from data, such as SuperPoint~\cite{detone_superpoint_2018} for joint keypoint detection and description, and SuperGlue~\cite{sarlin_superglue_2020} for learned correspondence matching, which substantially improve robustness under wide viewpoint changes, illumination differences and weakly textured surfaces where handcrafted descriptors tend to fail~\cite{jin_image_2021}. When two or more reference images can be matched to a session image, the 3D coordinates of the matched feature points can be triangulated from the known reference poses, and the session camera pose can be estimated through a perspective-n-point solver~\cite{jinyu_survey_2019}.

% Paragraph purpose: discuss limitations and challenges
A significant challenge for image-based alignment in construction and facility environments is the potential for drastic changes in appearance between the reference and test data. Lighting conditions, surface treatments, and the presence or absence of temporary objects can all alter the visual features of a scene. Derivative features such as vanishing lines or geometric edge features have been proposed as more robust alternatives, as they are less likely to be associated with temporary scene elements~\cite{zollmann_augmented_2014}. Additionally, the repetitive nature of many indoor environments can cause ambiguous matches, where features from different rooms or corridors appear visually identical.

% Paragraph purpose: discuss combined approaches
Several works have explored the combination of 2D and 3D alignment methods. Liu et al.~\cite{liu_instant_2018} initialise outdoor SLAM by estimating relative poses from localised panoramic images. Ventura et al.~\cite{ventura_global_2014} propose multi-view object detection combined with feature matching for global localisation. Brachmann and Rother~\cite{brachmann_expert_2019} extract CNN features and apply Expert Sample Consensus to handle scene outliers. However, most of these methods focus on either 2D or 3D data independently and do not combine both modalities in a unified pose estimation framework. The method presented in this chapter addresses this gap by jointly leveraging 2D image features, 3D point cloud features, and BIM geometry to produce a robust and accurate alignment.

\subsection{Dataset Updating}
\label{sec:4_rw_updating}

Beyond alignment, combining multi-temporal datasets into a single consistent representation remains an open challenge. The simplest case is coverage improvement in a static environment: Lachat~et~al.~\cite{lachat_combination_2016} supplement TLS data with close-range photogrammetry, while Bolkas~et~al.~\cite{bolkas_creating_2020} fuse drone and ground-level point clouds using a grid-based occlusion infilling strategy. Alteirac~et~al.~\cite{alteirac_registration_2021} combine simultaneous multi-sensor captures. These methods assume the scene is unchanged between acquisitions.

When the environment does change, the discrepancies must first be detected. Cloud-to-cloud distance computation is the dominant approach~\cite{mukupa_review_2017}, with voxel-based~\cite{hirt_change_2021}, learning-based~\cite{meyer_change_2022}, and comparative point-based variants such as M3C2~\cite{shirowzhan_comparative_2019} improving robustness. Nikoohemat~et~al.~\cite{nikoohemat_change_2018} specifically target indoor change detection from point clouds, though their scope is limited to identifying differences between exactly two epochs rather than integrating the new observations into the existing dataset. More broadly, most change detection methods share this limitation: they detect and report differences but do not produce an updated dataset.

Several works have addressed the progressive integration of new captures into an existing model. Gao~et~al.~\cite{gao_approach_2015} propose combining progressively captured point clouds, though they do not address the case where the scene itself has changed between captures. Pu\v{c}ko~et~al.~\cite{pucko_automated_2018} propose continuous updating of a 4D as-built point cloud through helmet-mounted low-precision scanners, which is conceptually close to the HoloLens-based capture used in this work, but their method merges partial scans without distinguishing between changed and unchanged regions. Partovi~et~al.~\cite{partovi_automatic_2021} focus on filling occlusions with photogrammetric data but do not handle temporal changes.

The key difficulty that none of these methods fully resolve is distinguishing, at the point level, between three categories: unchanged regions where the higher-quality original data should be preserved, genuinely changed regions where the new data must replace the old, and regions that the new sensor simply did not observe. The method presented in this chapter addresses this three-way distinction through a sequential coverage and visibility check that operates directly on the point cloud without requiring a semantic model.

% ─────────────────────────────────────────────────────────────────────────────

\begin{figure}[hb]
  \centering
  \includegraphics[width=1\textwidth]{Alignment_Method}
  \caption[Alignment pipeline overview]{Overview of the proposed two-step alignment and updating pipeline.}
  \label{fig:4_method_overview}
\end{figure}


\section{Methodology}
\label{sec:4_methodology}

The proposed alignment and updating pipeline consists of three sequential stages (Figure~\ref{fig:4_method_overview}). First, a global alignment narrows the search area for the pose estimation. Second, a local alignment computes the precise pose of the captured session using multiple complementary methods. Third, a dataset updating procedure integrates the newly captured data into the existing reference point cloud. Each stage is described in detail below.


% ── 4.3.1 ────────────────────────────────────────────────────────────────────
\subsection{Data Preprocessing and Session Management}
\label{sec:4_preprocessing}

% Paragraph purpose: explain the session concept and data organisation
When handling datasets from different sensors, time periods, and locations, a standardised data management framework is essential. In this work, each data capture is organised as a \textit{session}. A session contains all data captured during a single measurement campaign with a single device, including point clouds, meshes, images, and their associated metadata. Sessions can originate from any of the sensing technologies described in Chapter~\ref{ch:3_Data_Acquisition}: high-end TLS, mobile mapping systems, \gls{XR} devices, or smartphones.

% Paragraph purpose: explain the RDF-based metadata management
To enable interoperability between sessions captured by different sensors and at different times, the metadata of each session is serialised as an \gls{RDF} graph using the GEOMAPI framework~\cite{bassier_geomapi_2024}. The graph stores the geospatial metadata of each resource (position, orientation, bounding box, coordinate system, accuracy, timestamp, \dots) without duplicating the actual point cloud or image data. This lightweight representation allows rapid spatial queries across large collections of sessions without loading the underlying datasets.

% Paragraph purpose: explain reference data preprocessing
The reference data used for the alignment can consist of any combination of previously captured sessions: \gls{TLS} point clouds, \gls{iMMS} data, geolocated imagery, and \gls{BIM} models. Each reference resource is preprocessed as follows. Point clouds are converted to PCD format and a downsampled voxel octree together with \gls{FPFH} geometric features is extracted and stored in the RDF graph. Images have their \gls{ORB} features extracted during the \gls{SfM} processing and these features are stored per image in the graph. BIM geometries are sampled to point clouds and \gls{SUPER4PCS} features are extracted. In all cases, only the metadata and features are stored in the graph; the raw data remains at its original location.

% Paragraph purpose: explain XR session preprocessing
For each new XR measurement session, the captured images and meshes are sent to a processing server. The server extracts the relevant metadata and features and organises them in the session's RDF graph, analogous to the reference data preprocessing. The relative poses within the session are directly adopted from the device's local SLAM output, which is assumed to be sufficiently accurate for single-room environments~\cite{geyter_point_2022}.

% ── 4.3.2 ────────────────────────────────────────────────────────────────────
\subsection{Global Alignment}
\label{sec:4_global}

% Paragraph purpose: explain the purpose and mechanism of the global alignment
Due to the potentially large volume of reference data accumulated over time, it is computationally infeasible to match every new session against all available reference resources. This is why the first step of the alignment process is a coarse global positioning that restricts the search space for the subsequent fine alignment.

% Paragraph purpose: explain the positioning mechanism
The global position of the new session is obtained from the HTML Geolocation \gls{API}, which estimates the device's location based on available Wi-Fi, radio, and \gls{GNSS} signals~\cite{vermandere_two-step_2022}. In indoor environments, the resulting position estimate $[\mathbf{p}_\text{ses} \pm \sigma_\text{glob}]$ has an accuracy of approximately 20\,m, which is sufficient to narrow down the available reference data to the correct building, but not to achieve the metric-level accuracy required for data integration.

% Paragraph purpose: explain the reference data sub-selection
Given the coarse position estimate, relevant subsets of the reference point clouds and images are extracted. For images, the Euclidean distance between the focal point of each session image and each reference image is evaluated, retaining only reference images within a threshold distance $t_\text{dist}$ augmented by the positioning uncertainty $\sigma_\text{glob}$. For point clouds, any reference cloud whose bounding box intersects with the session's bounding box expanded by $\sigma_\text{glob}$ is retained:
\begin{equation}
\label{eq:4_global_selection}
\begin{split}
\mathcal{P}_\text{ref}' &= \left\{ P_\text{ref} \in \mathcal{P}_\text{ref} \;\middle|\; P_\text{ref} \cap \left[\mathbf{p}_{\text{ses},\min} - \sigma_\text{glob},\; \mathbf{p}_{\text{ses},\max} + \sigma_\text{glob}\right] \neq \emptyset \right\} \\
I_\text{ref}' &= \left\{ i_\text{ref} \in I_\text{ref} \;\middle|\; \exists\, i_\text{ses} \in I_\text{ses} : \|\mathbf{c}_\text{ref} - \mathbf{c}_\text{ses}\| \leq t_\text{dist} + \sigma_\text{glob} \right\}
\end{split}
\end{equation}
where:
\begin{itemize}
    \item $\mathcal{P}_\text{ref}$, $I_\text{ref}$ are the full reference point cloud and image collections,
    \item $\mathcal{P}_\text{ref}'$, $I_\text{ref}'$ are the selected subsets,
    \item $\mathbf{c}_\text{ref}$, $\mathbf{c}_\text{ses}$ are the camera centres of a reference and session image,
    \item $\sigma_\text{glob}$ is the global positioning uncertainty,
    \item $t_\text{dist}$ is the image distance threshold.
\end{itemize}
This selection is performed entirely on the RDF metadata graphs, making it extremely fast and avoiding the need to load any raw data at this stage.

% ── 4.3.3 ────────────────────────────────────────────────────────────────────
\subsection{Local Alignment}
\label{sec:4_local}

% Paragraph purpose: introduce the multi-method local alignment
Once the relevant reference data has been identified, the precise pose of the session is computed using multiple complementary alignment methods. Each method produces a candidate transformation together with quality metrics. These candidates are subsequently combined in a weighted pose vote to determine the final session pose. The methods are divided into 3D geometric alignment and 2D image-based alignment.

% ── 4.3.3.1 ──────────────────────────────────────────────────────────────────
\subsubsection{3D Alignment}
\label{sec:4_3d_alignment}

\begin{figure}[htbp]
  \centering
  \medskip
  \includegraphics[width=1\textwidth]{3D_matching_experiment}
  \caption[3D Alignment]{3D alignment of a session scan (middle) against a georeferenced reference scan (left), producing a precise alignment with the pre-aligned BIM model (right).}
  \label{fig:3D_Alignment_method}
\end{figure}

% Paragraph purpose: explain FPFH-based point cloud alignment
The first 3D alignment method matches the session's point cloud or mesh against reference point clouds using \gls{FPFH} feature correspondences (Figure~\ref{fig:3D_Alignment_method}). Given the 3D features of the reference and session datasets, a rigid body transformation is computed using the Fast Global Registration algorithm~\cite{zhou_fast_2016}. This method minimises the distances between feature correspondences while simultaneously rejecting outliers through a robust correspondence function $\rho$:
\begin{equation}
\label{eq:4_fpfh}
\mathbf{T}_\text{pc} = \argmin_{\mathbf{T}} \sum_{X_\text{ref},\,X_\text{ses}} \rho \left(\| X_\text{ref}- \mathbf{T}\, X_\text{ses}\|\right)
\end{equation}
where:
\begin{itemize}
    \item $X_\text{ref}$, $X_\text{ses}$ are the matched 3D feature points in the reference and session datasets,
    \item $\mathbf{T}$ is the candidate rigid body transformation,
    \item $\mathbf{T}_\text{pc}$ is the resulting point-cloud-based transformation,
    \item $\rho$ is the robust correspondence weighting function.
\end{itemize}
Only the precomputed feature graphs need to be transferred to the server, avoiding the computational cost of processing the full point clouds in real time.

% Paragraph purpose: explain BIM-based alignment
The second 3D alignment method targets the \gls{BIM} model using the \gls{SUPER4PCS} algorithm~\cite{mellado_super_2014}, described in Section~\ref{sec:4_related_3d}. Since BIM geometries lack the local surface detail that \gls{FPFH} descriptors rely on, \gls{SUPER4PCS} is used instead, as its ratio-based matching operates on sparse planar geometry without requiring rich surface normals. The sampled BIM point cloud and the session point cloud are both provided to \gls{SUPER4PCS}, which returns a candidate rigid body transformation together with the Largest Common Pointset (\gls{LCP}) score that quantifies the fraction of the input that participates in the alignment. This score, along with the inlier percentage and the spatial extent of the matched region, feeds into the pose vote described in Section~\ref{sec:4_pose_vote}.

% ── 4.3.3.2 ──────────────────────────────────────────────────────────────────
\subsubsection{2D Alignment}
\label{sec:4_2d_alignment}

% Paragraph purpose: introduce the three image-based alignment variants
Three image-based alignment methods are employed, each addressing a different configuration of available reference imagery. All methods are based on \gls{ORB} feature matching between session and reference images.

\begin{wrapfigure}[14]{l}{0.33\textwidth}
  \centering
  \includegraphics[width=0.33\textwidth]{2D_Alignment_overlapping}
  \caption[2D alignment overlapping]{Two overlapping reference images (red) matched to the session image (blue).}
  \label{fig:2D_Alignment_overlapping}
\end{wrapfigure}
% Paragraph purpose: explain the two overlapping references method
\paragraph{Two overlapping references.}
When a session image $i_\text{ses}$ can be matched with at least two overlapping co-registered reference images $i_{\text{ref},1}$ and $i_{\text{ref},2}$ with known transformations $\mathbf{T}_{\text{ref},1}$ and $\mathbf{T}_{\text{ref},2}$ (Figure~\ref{fig:2D_Alignment_overlapping}), the 3D coordinates of the matched features are already determined in the global coordinate system through the known reference poses. The session camera pose $\mathbf{T}_\text{ses}$ is then retrieved by solving the perspective-n-point problem using at least six 3D--2D correspondences. This method is analogous to conventional \gls{SfM} processing and produces the most reliable image-based alignment when dense, overlapping reference imagery is available.

\begin{wrapfigure}[11]{r}{0.33\textwidth}
  \centering
  \vspace{-20pt}
  \includegraphics[width=0.33\textwidth]{2D_Alignment_separate}
  \caption[2D alignment separate]{Two separate reference images (red) matched to the session image (blue).}
  \label{fig:2D_Alignment_separate}
\end{wrapfigure}
% Paragraph purpose: explain the two separate references method
\paragraph{Two separate references.}
When two non-overlapping reference images $i_{\text{ref},1}$ and $i_{\text{ref},2}$ with known poses $\mathbf{T}_{\text{ref},1}$ and $\mathbf{T}_{\text{ref},2}$ can each be matched to the session image $i_\text{ses}$ (Figure~\ref{fig:2D_Alignment_separate}), each match recovers a relative rotation and translation direction through the essential matrix decomposition. Since a single image pair provides no scale, only the direction from the reference camera to the session camera is known, not the distance. The session pose $\mathbf{T}_\text{ses}$ is therefore recovered by intersecting the two direction rays originating from the known reference positions:
\begin{equation}
\label{eq:4_separate_ref}
\mathbf{T}_\text{ses} = \argmin_{\mathbf{T}} \sum_{k \in \{1,2\}} \left\| \mathbf{T} - \mathbf{T}_{\text{ref},k} \cdot \hat{\mathbf{T}}_{\text{rel},k} \right\|
\end{equation}
where $\hat{\mathbf{T}}_{\text{rel},k}$ is the unit-scale relative transformation between $i_\text{ses}$ and $i_{\text{ref},k}$. The accuracy depends on the angular separation between the two reference viewpoints: near-perpendicular configurations yield a well-defined intersection, while near-parallel ones amplify directional errors into large positional uncertainty.

\begin{wrapfigure}[15]{l}{0.33\textwidth}
  \centering
  \vspace{-5pt}
  \includegraphics[width=0.33\textwidth]{2D_Alignment_geometry}
  \caption[2D alignment geometry]{Single reference image with associated 3D geometry (blue) matched to the session image (red).}
  \label{fig:2D_Alignment_geometry}
\end{wrapfigure}
% Paragraph purpose: explain the single reference with geometry method
\paragraph{One reference image with geometry.}
When only a single reference image match is available but 3D geometry is present in either the session or the reference data (Figure~\ref{fig:2D_Alignment_geometry}), the scale of the transformation is recovered through raycasting (Figure~\ref{fig:Raycasting_method}). A set of rays is constructed from the camera focal point through the matched 2D feature points, and the intersections with the available 3D geometry yield the 3D coordinates of the correspondences:
\begin{equation}
\label{eq:4_raycast}
X=\left\{ \mathbf{p} \in P \;\middle|\; \exists\, d(\mathbf{o},\mathbf{c}) \in D: \mathbf{p}= d(\mathbf{o},\mathbf{c}) \cap P\right\}
\end{equation}
where:
\begin{itemize}
    \item $D$ is the set of rays constructed from the focal point $\mathbf{o}$ through the matched pixels $\mathbf{c}$,
    \item $P$ is the available 3D geometry,
    \item $X$ is the resulting set of recovered 3D correspondence points.
\end{itemize}
This method is particularly useful when only sparse imagery is available but the session contains depth information from a ToF sensor.

\begin{figure}[htbp]
  \centering
  \medskip
  \includegraphics[width=1\textwidth]{Good_raycasting_experiment}
  \caption[Raycasting Alignment]{Single-reference raycasting alignment: the session image (right) is matched against a reference image with associated geometry (middle and left), recovering the transformation scale through ray intersections.}
  \label{fig:Raycasting_method}
\end{figure}

% ── 4.3.3.3 ──────────────────────────────────────────────────────────────────
\subsubsection{Final Pose Estimation}
\label{sec:4_pose_vote}

% Paragraph purpose: explain the weighted voting mechanism
Each alignment method produces a candidate transformation together with quality parameters. To combine these into a single robust pose estimate, a weighted pose vote is computed over all candidate transformations. Since rigid body transformations comprise a translation and a rotation, and rotations do not form a vector space, the two components are averaged separately.

The final position $\mathbf{t}_\text{ses}$ is computed as a standard weighted average over the candidate translations:
\begin{equation}
\label{eq:4_pose_vote_pos}
\mathbf{t}_\text{ses} = \frac{\sum_{i=1}^{n} w_i \, \mathbf{t}_i}{\sum_{i=1}^{n} w_i}
\end{equation}

The final orientation is computed through weighted quaternion averaging \cite{cheng_averaging_2007}. Each candidate rotation is represented as a unit quaternion $\mathbf{q}_i$, and the weighted average quaternion $\mathbf{q}_\text{ses}$ is obtained as the eigenvector corresponding to the largest eigenvalue of the weighted quaternion matrix:
\begin{equation}
\label{eq:4_pose_vote_rot}
\mathbf{q}_\text{ses} = \text{eig}_\text{max}\!\left(\sum_{i=1}^{n} w_i \, \mathbf{q}_i \mathbf{q}_i^\top\right)
\end{equation}
where:
\begin{itemize}
    \item $\mathbf{t}_i$, $\mathbf{q}_i$ are the translation and unit quaternion of candidate $i$,
    \item $w_i$ is the weight of candidate $i$,
    \item $n$ is the number of candidate transformations.
\end{itemize}

% Paragraph purpose: explain the weight computation
For image-based transformations, the weight $w_i$ is derived from the reprojection error, the number of feature matches, the inlier percentage, and the spatial extent of the inlier correspondences. For point cloud-based transformations, the weight is based on the \gls{RMSE} of the matched correspondences, the number of matches, and the spatial extent of the match region. Additionally, a method-specific base weight reflects the empirical reliability of each approach: \gls{SUPER4PCS} alignment receives the highest base weight (1.0), followed by \gls{FPFH}-based registration and \gls{SfM} with overlapping references (both 0.8), two separate reference images (0.5), and single-reference raycasting (0.4).

% Paragraph purpose: explain the output
The result of the pose voting is both a position and orientation estimate for the session and a confidence measure that reflects the degree of agreement among the contributing methods. This confidence measure is valuable for downstream quality control, as it indicates how reliably the session has been positioned within the reference coordinate system.
% ── 4.3.4 ────────────────────────────────────────────────────────────────────
\subsection{New Data Combination}
\label{sec:4_combination}

% Paragraph purpose: motivate the need for point cloud updating
Once the newly captured session is aligned to the reference coordinate system, the new observations can be used to update the existing reference point cloud. The goal is to extend the lifespan of the high-quality reference data by replacing only the portions that have changed, while preserving the more accurate existing data wherever possible. This requires determining, at the point level, which observations in the reference dataset are still current and which have become outdated~\cite{vermandere_automatic_2022, bassier_geomapi_2024}.



% Paragraph purpose: introduce the three-phase updating process
The updating process consists of three sequential phases: the removal of out-of-date points from the reference dataset, a visibility check, and the addition of new points from the session dataset.


% ── 4.3.4.1 ──────────────────────────────────────────────────────────────────
\subsubsection{Out-of-Date Point Removal}
\label{sec:4_removal}

% Paragraph purpose: explain the sub-selection step
\begin{wrapfigure}[9]{r}{0.25\textwidth}
  \centering
  \vspace{-30pt}
  \includegraphics[]{Updating_method_removal}
  \caption[Point removal]{point removal}
  \label{fig:Updating_method_removal}
\end{wrapfigure}
To prevent false positives in regions not covered by the new session, only the portion of the reference point cloud that falls within the spatial extent of the new data is considered for updating. The convex hull of the session's point cloud is used to define this region, providing a tighter bound than an axis-aligned bounding box and reducing the number of incorrectly flagged points.

% Paragraph purpose: explain the coverage check
Within this sub-selection, each reference point is subjected to two sequential checks: a coverage check and a visibility check. The coverage check determines whether a reference point $\mathbf{p}_i$ has a corresponding observation in the new dataset by evaluating the nearest-neighbour distance to the session point cloud $Q$:

% Paragraph purpose: explain the sub-selection step
\begin{equation}
\label{eq:4_coverage}
P' = \left \{ \mathbf{p}_i \;\middle|\; \forall \mathbf{p}_i \in P : \min_{\mathbf{q}_j \in Q}\|\mathbf{p}_i-\mathbf{q}_j\| \geq t_\text{dist}\right\}
\end{equation}
where:
\begin{itemize}
    \item $P$ is the reference point cloud within the convex hull of the session,
    \item $Q$ is the session point cloud,
    \item $\mathbf{p}_i$, $\mathbf{q}_j$ are reference and session points,
    \item $t_\text{dist}$ is the distance threshold,
    \item $P'$ is the subset of uncovered reference points.
\end{itemize}
Points for which no nearby session observation exists are considered uncovered and form the subset $P'$.



% Paragraph purpose: explain the visibility check
\subsubsection{Visibility Check}
\label{sec:4_visibility}
\begin{wrapfigure}[11]{r}{0.25\textwidth}
  \centering
  \vspace{-10pt}
  \includegraphics[]{Updating_method_visibility}
  \caption[Point visibility]{point visibility}
  \label{fig:Updating_method_visibility}
\end{wrapfigure}
However, not all uncovered points are necessarily outdated. Some may simply have been occluded during the new session. For instance, a point on the other side of a wall that was not within the sensor's field of view. To distinguish between genuinely outdated points and occluded ones, a visibility check is performed. For each uncovered point $\mathbf{p}_i \in P'$, the normals of the $n_\text{nb}$ closest points $\mathbf{q}_j$ in the session dataset are evaluated. If the majority of these normals face away from $\mathbf{p}_i$, the point is considered not visible and is retained in the reference dataset. Otherwise, it is flagged as visible and removed:

\begin{equation}
\label{eq:4_visibility}
\begin{gathered}
Q_i = \left\{\mathbf{q}_j \;\middle|\; \mathbf{q}_j \in Q \text{ among the } n_\text{nb} \text{ nearest to } \mathbf{p}_i \right\} \\
Q_i' = \left\{\mathbf{q}_j \;\middle|\; \forall \mathbf{q}_j \in Q_i: \overrightarrow{\mathbf{p}_i\mathbf{q}_j} \cdot \vec{\mathbf{n}}(\mathbf{q}_j) \geq t_\text{norm}\right\} \\
P'' = \left \{ \mathbf{p}_i \;\middle|\; \forall \mathbf{p}_i \in P': \frac{|Q_i'|}{|Q_i|} \leq t_\text{vis}\right\}
\end{gathered}
\end{equation}
where:
\begin{itemize}
    \item $Q_i$ is the set of $n_\text{nb}$ nearest session points around $\mathbf{p}_i$,
    \item $Q_i'$ is the subset of $Q_i$ whose normals $\vec{\mathbf{n}}(\mathbf{q}_j)$ face away from $\mathbf{p}_i$,
    \item $t_\text{norm}$ is the normal alignment threshold,
    \item $t_\text{vis}$ is the visibility percentage threshold,
    \item $P''$ is the subset of uncovered points retained as occluded.
\end{itemize}
Points that are both uncovered and visible are removed from the reference dataset.


% ── 4.3.4.2 ──────────────────────────────────────────────────────────────────
\subsubsection{New Point Introduction}
\label{sec:4_addition}
\begin{wrapfigure}[]{r}{0.25\textwidth}
  \centering
  \includegraphics[]{Updating_method_addition}
  \caption[Point introduction]{point introduction}
  \label{fig:Updating_method_addition}
\end{wrapfigure}
% Paragraph purpose: explain the addition of new points
After removing outdated points, new observations from the session dataset are added to fill the resulting gaps and to represent newly appeared objects. Since the reference point cloud is assumed to have higher accuracy than the session data, only session points that have no existing counterpart in the (updated) reference dataset are added. This is determined by an inverse distance query: session points whose nearest neighbour in the remaining reference dataset exceeds the threshold $t_\text{dist}$ are appended:
\begin{equation}
\label{eq:4_addition}
Q' = \left \{ \mathbf{q}_j \;\middle|\; \forall \mathbf{q}_j \in Q : \min_{\mathbf{p}_i \in \{ P \setminus P'' \}}\|\mathbf{q}_j-\mathbf{p}_i\| \geq t_\text{dist}\right\}
\end{equation}
where:
\begin{itemize}
    \item $Q$ is the session point cloud and $\mathbf{q}_j$ a session point,
    \item $\{ P \setminus P'' \}$ is the reference point cloud after removal of the outdated points,
    \item $t_\text{dist}$ is the distance threshold,
    \item $Q'$ is the set of new session points appended to the reference.
\end{itemize}

% Paragraph purpose: explain the final combined dataset
The final updated dataset consists of three parts: the original reference points outside the convex hull of the session (unaffected by the updating), the reference points inside the convex hull that passed both checks (still current), and the newly added session points. This composition ensures that the more accurate reference data is preserved wherever possible, while outdated regions are replaced by the most recent observations. Colour information for the newly added points, which may not be present in the session mesh, can be recovered by projecting the localised session images onto the new points and assigning the corresponding pixel values.

% ─────────────────────────────────────────────────────────────────────────────
\section{Experiments}
\label{sec:4_experiments}

The alignment and updating methodology is evaluated on two complementary data sources. First, the three real-world environments introduced in Section~\ref{sec:3_captured_scenes} provide 45 measurement sessions spanning four sensor modalities under realistic multi-temporal conditions. Second, the V-Scan synthetic dataset (Section~\ref{sec:3_synthetic}) provides controlled ground-truth poses against which alignment accuracy can be measured precisely.

% ── 4.4.1 ────────────────────────────────────────────────────────────────────
\subsection{Alignment on Real-World Data}
\label{sec:4_alignment_experiments}

Table~\ref{tab:4_alignment_results} reports the positional and rotational accuracy of each individual alignment method and the combined weighted estimate for all three real-world test cases. The accuracy is measured against manually aligned ground-truth poses. Across all sites, the combined pose estimation achieves an average positional RMSE of 0.058\,m, an average rotational RMSE of 0.337°, and an average confidence of approximately 50\%.

\begin{table}[hb]
    \centering
    \caption[Real-world alignment results]{Positional and rotational accuracy and confidence of each alignment method across the three real-world test cases. Case~1: laboratory rooms, Case~2: living lab, Case~3: Bruges renovation.}
    \label{tab:4_alignment_results}
    \renewcommand*{\arraystretch}{1.2}
    \setlength{\tabcolsep}{4pt}
    \begin{tabularx}{\textwidth}{r|XXXXX|X}
        \hline
        Method
            & Linked\newline Ref
            & Separate\newline Ref
            & Ray-\newline casting
            & FPFH
            & SUPER\newline 4PCS
            & Combined \\
        Est.\ time
            & 1--10\newline img/s
            & 1--10\newline img/s
            & 0.2--2\newline img/s
            & 5s--2\newline min/pcd
            & 20s--5\newline min/pcd
            & 20s--15\newline min/ses. \\
        Weight
            & 0.8 & 0.5 & 0.4 & 0.8 & 1 & \\ \hline
        %
        \multicolumn{7}{l}{\textit{Case 1}} \\ [-0.5ex]
        \quad Pos.\ RMSE (m)  & 0.145 & 0.189 & 0.021 & 0.159 & 0.061 & 0.065 \\
        \quad Rot.\ RMSE (°)  & 0.960 & 0.636 & 0.189 & 1.507 & 0.463 & 0.538 \\
        \quad Conf.\ (\%)     & 30.1  & 29.9  & 52.2  & 33.9  & 77.3  & 47.0  \\ \hline
        %
        \multicolumn{7}{l}{\textit{Case 2}} \\ [-0.5ex]
        \quad Pos.\ RMSE (m)  & 0.183 & 0.133 & 0.277 & 0.059 & 0.071 & 0.035 \\
        \quad Rot.\ RMSE (°)  & 0.993 & 1.207 & 0.018 & 0.156 & 0.028 & 0.074 \\
        \quad Conf.\ (\%)     & 25.1  & 17.0  & 19.2  & 80.0  & 90.0  & 54.4  \\ \hline
        %
        \multicolumn{7}{l}{\textit{Case 3}} \\ [-0.5ex]
        \quad Pos.\ RMSE (m)  & 0.149 & 0.103 & 0.106 & 0.067 & 0.078 & 0.074 \\
        \quad Rot.\ RMSE (°)  & 0.852 & 1.073 & 0.713 & 0.966 & 0.171 & 0.399 \\
        \quad Conf.\ (\%)     & 20.2  & 21.9  & 31.4  & 67.0  & 75.0  & 48.1  \\ \hline
        %
        \multicolumn{7}{l}{\textit{Average}} \\ [-0.5ex]
        \quad Pos.\ RMSE (m)  & 0.159 & 0.141 & 0.135 & 0.095 & 0.070 & 0.058 \\
        \quad Rot.\ RMSE (°)  & 0.935 & 0.972 & 0.307 & 0.886 & 0.221 & 0.337 \\
        \quad Conf.\ (\%)     & 25.1  & 22.9  & 34.3  & 60.3  & 80.8  & 49.8  \\ \hline
    \end{tabularx}
\end{table}

Each of the three real-world sites poses distinct alignment challenges. The laboratory rooms test multi-temporal registration under object-level layout changes within a stable structural envelope. The living lab tests robustness when both geometry and surface appearance change drastically between epochs. The Bruges renovation tests indoor alignment in a cluttered residential environment with significant temporary storage of construction materials. The alignment was performed using the five methods described in Section~\ref{sec:4_local}, with parameters determined empirically across all sessions. The global distance threshold $t_\text{dist}$ was set to 10\,m. The correspondence function $\rho$ was set to 50 pixels for image matches and to 1.5 times the voxel size (0.05\,m) for point cloud matches.

The geometric methods (FPFH and SUPER4PCS) yield the highest individual confidences, averaging 60\% and 81\% respectively, with SUPER4PCS achieving the lowest average positional RMSE (0.070\,m) when a BIM reference is available. The image-based methods achieve lower individual confidences (20--35\%) but contribute a larger number of candidate matches, improving overall robustness through the pose voting mechanism. The combined estimate consistently outperforms any individual method in positional accuracy, confirming that the complementary use of 2D and 3D methods is essential for robust indoor pose estimation across changing environments.

\begin{figure}[htbp]
  \centering
  \includegraphics[width=1\textwidth]{Good_2D_matching_experiment}
  \caption[2D Matching Experiment]{Example of a successful 2D alignment in the living lab. Two session images (left and middle) are matched against a reference image (right).}
  \label{fig:2D_Alignment_exp}
\end{figure}

% ── 4.4.2 ────────────────────────────────────────────────────────────────────
\subsection{Alignment on Synthetic Data}
\label{sec:4_alignment_synthetic}

While the real-world experiments demonstrate practical robustness, they rely on manually established ground-truth poses whose own accuracy is limited by the registration procedure. The V-Scan dataset (Section~\ref{sec:3_virtual_scanner}) provides exact ground-truth scanner poses, enabling a more precise evaluation of alignment accuracy under controlled conditions.

The experiment aligns each scan variation against two types of reference data: the \textit{reference} scan (the main furnished scan of the same room) and the \textit{empty} scan (the ground-truth empty scene). Three variation types are evaluated, corresponding to the scan variation system described in Section~\ref{sec:3_scan_management}:

\begin{itemize}
    \item \textbf{Variation~0} --- furniture change: the room structure is preserved but the furniture layout is regenerated, simulating a temporal change in object configuration.
    \item \textbf{Variation~1} --- position change: the scanner is repositioned within the refurnished room, simulating a different measurement setup.
    \item \textbf{Variation~2} --- scanner change: a lower-resolution scanner preset is used along with a new scanner position, simulating a heterogeneous multi-sensor scenario.
\end{itemize}

\paragraph{Full scan alignment.}

Table~\ref{tab:4_vscan_full} reports the alignment results when the variation scans cover the full room (Figure~\ref{fig:2D_Alignment_exp_synthetic_full}). The alignment is evaluated across all 200 V-Scan scenes and the success rate, average RMSE, positional error, and angular error are reported for successful registrations.

\begin{figure}[hb]
  \centering
  \includegraphics[width=1\textwidth]{Synthetic_alignment_results_full}
  \caption[Synthetic Alignment Results (Full Scans)]{Example of a successful full scan alignment in Variation~2 (Furniture Change). The variation scan (red) is aligned against the furnished reference scan (blue, top) and the empty scene scan (blue, bottom).}
  \label{fig:2D_Alignment_exp_synthetic_full}
\end{figure}

\begin{table}[hb]
    \centering
    \caption[V-Scan full scan alignment]{Alignment results on V-Scan full scans. Each variation is aligned against the furnished reference scan and the empty scene scan. Only successful registrations are included in the accuracy statistics.}
    \label{tab:4_vscan_full}
    \begin{tabular}{l l r r r}
        \toprule
        & & \textbf{Var.~0} & \textbf{Var.~1} & \textbf{Var.~2} \\
        & & (furniture) & (position) & (scanner) \\
        \midrule
        \multirow{4}{*}{Reference}
            & Success (\%)         & 88.0  & 87.0  & 61.0  \\
            & Avg.\ RMSE (m)       & 0.023 & 0.028 & 0.050 \\
            & Avg.\ pos.\ error (m) & 0.020 & 0.025 & 0.009 \\
            & Avg.\ ang.\ error (°) & 0.014 & 0.014 & 0.024 \\
        \midrule
        \multirow{4}{*}{Empty}
            & Success (\%)         & 96.0  & 87.0  & 55.0  \\
            & Avg.\ RMSE (m)       & 0.022 & 0.027 & 0.050 \\
            & Avg.\ pos.\ error (m) & 0.027 & 0.026 & 0.017 \\
            & Avg.\ ang.\ error (°) & 0.024 & 0.012 & 0.015 \\
        \bottomrule
    \end{tabular}
\end{table}

Furniture changes (Variation~0) and position changes (Variation~1) are handled reliably, with success rates of 87--96\% and sub-3\,cm RMSE. Aligning against the empty scene yields even higher success for Variation~0 (96\% vs.\ 88\%), which is expected since the structural geometry --- walls, floor, ceiling --- remains identical and provides unambiguous correspondences once the furniture clutter is removed. The scanner change (Variation~2) is the most challenging condition, reducing the success rate to 55--61\% and increasing the RMSE to 0.050\,m. This is consistent with the degraded geometric detail produced by the lower-resolution scanner preset, which provides fewer distinctive features for the registration algorithm.

\paragraph{Partial scan alignment.}

In practice, new observations often cover only a portion of the existing scene. To evaluate this scenario, each variation scan is cropped into multiple overlapping partial views, and the alignment is repeated on each crop (Figure~\ref{fig:2D_Alignment_exp_synthetic_partial}). Table~\ref{tab:4_vscan_partial} reports the results averaged across all partial crops.

\begin{figure}[htbp]
  \centering
  \includegraphics[width=1\textwidth]{Synthetic_alignment_results_partial}
  \caption[Synthetic Alignment Results (Partial Scans)]{Example of a successful partial scan alignment in Variation~2 (Furniture Change). A random crop of the variation scan (red) is aligned against the furnished reference scan (blue, top) and the empty scene scan (blue, bottom).}
  \label{fig:2D_Alignment_exp_synthetic_partial}
\end{figure}

\begin{table}[htbp]
    \centering
    \caption[V-Scan partial scan alignment]{Alignment results on V-Scan partial scans (averaged across multiple crops per scene). Each variation is aligned against the furnished reference and the empty scene.}
    \label{tab:4_vscan_partial}
    \begin{tabular}{l l r r r}
        \toprule
        & & \textbf{Var.~0} & \textbf{Var.~1} & \textbf{Var.~2} \\
        & & (furniture) & (position) & (scanner) \\
        \midrule
        \multirow{4}{*}{Reference}
            & Success (\%)         & 66.7  & 68.0  & 22.0  \\
            & Avg.\ RMSE (m)       & 0.028 & 0.026 & 0.052 \\
            & Avg.\ pos.\ error (m) & 0.021 & 0.023 & 0.024 \\
            & Avg.\ ang.\ error (°) & 0.044 & 0.031 & 0.030 \\
        \midrule
        \multirow{4}{*}{Empty}
            & Success (\%)         & 82.7  & 72.7  & 22.0  \\
            & Avg.\ RMSE (m)       & 0.025 & 0.023 & 0.052 \\
            & Avg.\ pos.\ error (m) & 0.016 & 0.012 & 0.012 \\
            & Avg.\ ang.\ error (°) & 0.022 & 0.026 & 0.040 \\
        \bottomrule
    \end{tabular}
\end{table}

The partial scan scenario reduces the success rate by approximately 15--20 percentage points across all variations compared to full scans, reflecting the reduced geometric overlap available for correspondence matching. The RMSE of successful registrations, however, remains comparable (0.023--0.052\,m), indicating that when the alignment succeeds, it achieves similar accuracy regardless of coverage. The scanner change condition (Variation~2) drops to a 22\% success rate for partial scans, confirming that the combination of reduced resolution and reduced spatial coverage presents the most challenging registration scenario.

Aligning against the empty scene again outperforms the furnished reference for Variations~0 and~1, with a notably higher success rate for partial furniture-change scans (82.7\% vs.\ 66.7\%). This advantage diminishes for the scanner change, where both references yield identical success rates. The consistent advantage of the empty-scene reference suggests that the structural geometry of the room provides more reliable correspondences than the furniture, which may have changed or may introduce ambiguous matches between similar objects.

% ── 4.4.3 ────────────────────────────────────────────────────────────────────
\subsection{Combination Experiments}
\label{sec:4_combination_experiments}

The dataset updating procedure is evaluated in two settings: a real-world multi-sensor scenario that tests the method under practical conditions, and a controlled synthetic scenario on the V-Scan dataset that enables quantitative assessment of coverage improvement against known ground truth.

\subsubsection{Real-World Updating}

\begin{figure}[htbp]
  \centering
  \medskip
  \includegraphics[width=1\textwidth]{Updating_experiment}
  \caption[Updating Experiment]{The reference scan (top left) is updated with the new session scan (top right) to produce the combined dataset (bottom). The reference points are coloured by their status: removed (red) and updated (blue).}
  \label{fig:Updating_Experiment}
\end{figure}

The real-world evaluation uses a subset of the laboratory room data (Lab~2, Asphalt). Two captures of the same space are combined (Figure~\ref{fig:Updating_Experiment}): the reference dataset acquired with the NavVis VLX (743,618 points, 0.005\,m density) and a new session acquired with the HoloLens~2 (13,107 points, 0.05\,m density). Between captures, objects were added, removed, and repositioned, creating the mixed-change scenario described in Section~\ref{sec:4_rw_updating}.

The convex hull filtering reduced the relevant portion of the reference dataset by 95.5\%, confirming the effectiveness of the spatial sub-selection. During the removal phase, 36.2\% of the relevant reference points were not covered by the new session and were subjected to the visibility check. Of these, 56.6\% were determined to be visible and therefore outdated, resulting in the removal of 15.7\% of the relevant reference points. In the addition phase, 26.4\% of the session points had no existing counterpart and were added.

The final combined dataset consists of 54.4\% preserved reference points, 13.4\% retained but inconclusive occluded points, and 32.2\% newly added session points. The complete updating algorithm executed in 1.47\,s on a standard notebook. The result preserves the high accuracy of the VLX data in unchanged regions while incorporating HoloLens~2 observations where updates have occurred.

\subsubsection{Synthetic Updating}

The Ground truth available in the V-Scan dataset provides an objective reference against which coverage can be measured. For each of the five V-Scan environment types, the main furnished scan is combined with a furniture-and-position-change variation (Variation~1) using the updating procedure (Figure~\ref{fig:Synthetic_Combination_Results}). The reference and variation scans are both subsampled to a 2\,cm voxel resolution with a distance threshold of 2\,cm. Coverage is measured as the percentage of the ground-truth empty scene that is represented within the distance threshold by either the variation scan alone or the combined pipeline output.

\begin{figure}[hb]
  \centering
  \includegraphics[width=1\textwidth]{Synthetic_combination_results}
  \caption[Synthetic Combination Results (Full Scans)]{Example of a successful combination of a variation scan (red) with the reference scan (blue) to produce the combined dataset.}
  \label{fig:Synthetic_Combination_Results}
\end{figure}

Table~\ref{tab:4_vscan_combination} reports the results per environment type. Across all environments, the updating procedure preserves on average 92\% of the reference points and replaces approximately 8\% with observations from the variation scan. The replaced points correspond to regions where the furniture configuration changed between the two scans. The coverage of the empty scene improves from 89\% for the variation scan alone to over 97\% after the combination, yielding an average coverage gain of 8.4 percentage points.

\begin{table}[hb]
    \centering
    \caption[V-Scan combination results]{Updating results on V-Scan scenes per environment type. Coverage is measured against the ground-truth empty scene. All scans are subsampled to 2\,cm voxel size with a 2\,cm distance threshold.}
    \label{tab:4_vscan_combination}
    \setlength{\tabcolsep}{6pt}
    \resizebox{\textwidth}{!}{%
    \begin{tabular}{l c c c c c}
        \toprule
        \textbf{Environment}
            & \textbf{Points Kept}
            & \textbf{Points Updated}
            & \textbf{Var.\ Coverage}
            & \textbf{Pipeline Cov.}
            & \textbf{Coverage Gain} \\
        & (\%) & (\%) & of Empty (\%) & of Empty (\%) & (\%) \\
        \midrule
        Industrial   & 89.47 & 10.53 & 83.34 & 96.23 & $+$12.90 \\
        Office       & 90.71 &  9.29 & 87.19 & 96.70 & $+$9.51  \\
        Bedroom      & 93.61 &  6.39 & 93.81 & 98.52 & $+$4.70  \\
        Dining Room  & 92.25 &  7.75 & 88.27 & 97.57 & $+$9.29  \\
        Living Room  & 94.31 &  5.69 & 92.72 & 98.38 & $+$5.66  \\
        \midrule
        \textbf{Average} & \textbf{92.07} & \textbf{7.93} & \textbf{89.07} & \textbf{97.48} & \textbf{$+$8.41} \\
        \bottomrule
    \end{tabular}}
\end{table}

The results reveal a clear relationship between scene complexity and updating behaviour. The industrial environment, which contains the densest furniture arrangement, shows the highest proportion of updated points (10.5\%) and the largest coverage gain (+12.9 pp), reflecting the greater number of regions where furniture occlusions differ between the two scans. Conversely, the bedroom and living room environments, which contain fewer and more sparsely distributed objects, show lower update rates (5.7--6.4\%) and smaller coverage gains (+4.7--5.7 pp), as fewer regions are affected by the furniture change.

The fact that the combined coverage consistently exceeds 96\% of the empty scene confirms that the updating procedure effectively recovers structural surfaces that were occluded by different objects in each individual scan. The remaining 2.5\% of uncovered empty-scene geometry corresponds to regions that are occluded in both scans, such as floor areas beneath large furniture items that were not displaced between variations, or wall surfaces behind objects that remained in place.

% ─────────────────────────────────────────────────────────────────────────────
\section{Conclusion}
\label{sec:4_conclusion}

This chapter presented a two-step alignment framework that positions heterogeneous sensor data within a pre-existing reference frame through a coarse global positioning step followed by a fine registration step, unifying five complementary 2D and 3D methods in a weighted pose vote, together with a mechanism for updating the reference point cloud with new observations.

Across three real-world test cases the combined estimate achieves an average positional accuracy of 0.058\,m, meeting or exceeding the resolution of the lowest-quality sensors used. No individual method is reliable across all conditions, but the voting mechanism compensates for their respective weaknesses, with geometric methods providing high accuracy and image-based methods broad candidate coverage. The synthetic V-Scan experiments confirm this under exact ground truth: success rates of 87--96\% for furniture and position changes with sub-3\,cm RMSE, degrading gracefully under partial coverage and more sharply under scanner changes, which combined with reduced overlap is the hardest condition. A consistent finding across both is that structural geometry yields more reliable correspondences than furnished scenes, as the empty-scene reference outperforms the furnished one throughout.

The updating methodology enables incremental maintenance without full re-documentation: the two-phase removal and addition procedure distinguishes outdated, occluded, and current observations, preserving high-quality reference data while adding new observations only where needed. On V-Scan, combining two scans with different furniture configurations recovers on average 97.5\% of the ground-truth empty scene, compared to 89\% from a single scan.

This establishes the aligned, up-to-date point cloud that serves as input for all subsequent stages, which assume spatially consistent data representing the current state of the environment, even when the input originates from heterogeneous sensors captured at different times.


%%%%%%%%%%%%%%%%%%%%%%%%%%%%%%%%%%%%%%%%%%%%%%%%%%
\cleardoublepage